% When the encoding lies: how disjoint_EL turned a true LN lemma false.
% Standalone TeX article documenting the chain of cascading design
% choices behind the claim_3_25 "disproof". Compiles via
%   latexmk -pdflua 01_disjoint_EL_cascade.tex
% from this folder. (Requires lualatex because the embedded Lean code
% blocks contain UTF-8 mathematical symbols.)
% !TEX program = lualatex

\documentclass[11pt,a4paper]{article}

% Minimal self-contained preamble. Cannot use the leanification
% preamble because it loads inputenc=latin1, which collides with the
% UTF-8 Lean code blocks below. Only the macros actually referenced in
% this article are defined.
\usepackage[utf8]{inputenc}
\usepackage[T1]{fontenc}
\usepackage{amsmath,amssymb,amsthm}
\usepackage{stmaryrd}
\usepackage{xcolor}
\usepackage{tikz}
\usetikzlibrary{arrows.meta,decorations.markings}
\usepackage{hyperref}
\hypersetup{colorlinks=true,linkcolor=blue,urlcolor=blue}
\usepackage{listings}

% LN graph-notation macros (subset of leanification/preamble.tex).
\providecommand{\arrhead}{{Latex}}
\providecommand{\arrtail}{{}}
\providecommand{\arrstar}{Rays[n=6]}
\providecommand{\tuh}[1][]{\mathrel{\tikz [baseline=-0.25ex,\arrtail-\arrhead, #1] \draw [#1] (0pt,0.5ex) -- (1.3em,0.5ex);}}
\providecommand{\hut}[1][]{\mathrel{\tikz [baseline=-0.25ex,\arrhead-\arrtail, #1] \draw [#1] (0pt,0.5ex) -- (1.3em,0.5ex);}}
\providecommand{\huh}[1][]{\mathrel{\tikz [baseline=-0.25ex,\arrhead-\arrhead, #1] \draw [#1] (0pt,0.5ex) -- (1.3em,0.5ex);}}
\providecommand{\sus}[1][]{\mathrel{\tikz [baseline=-0.25ex,{Rays[n=6]}-{Rays[n=6]}, #1] \draw [#1] (0pt,0.5ex) -- (1.3em,0.5ex);}}
\providecommand{\hus}[1][]{\mathrel{\tikz [baseline=-0.25ex,\arrhead-{Rays[n=6]}, #1] \draw [#1] (0pt,0.5ex) -- (1.3em,0.5ex);}}
\providecommand{\suh}[1][]{\mathrel{\tikz [baseline=-0.25ex,{Rays[n=6]}-\arrhead, #1] \draw [#1] (0pt,0.5ex) -- (1.3em,0.5ex);}}

% Set / operator macros.
\providecommand{\sm}{\setminus}
\providecommand{\ins}{\subseteq}
\providecommand{\lp}{\left ( }
\providecommand{\rp}{\right ) }
\providecommand{\lC}{\left \{ }
\providecommand{\rC}{\right \} }
\providecommand{\given}{\mathbin{|}}

% Independence / separation notation.
\providecommand{\Perp}{\operatorname{\perp}}
\providecommand{\nPerp}{\operatorname{\not\perp}}
\providecommand{\iPerp}{\Perp^i}
\providecommand{\niPerp}{\nPerp^i}
\providecommand{\isPerp}{\Perp^{i\sigma}}
\providecommand{\nisPerp}{\nPerp^{i\sigma}}

% Graph predicates.
\providecommand{\Anc}{\mathrm{Anc}}
\providecommand{\Desc}{\mathrm{Desc}}
\providecommand{\Sc}{\mathrm{Sc}}

% A bare theorem env so we can render the lemma statement.
\newtheorem*{Lem}{Lemma}

% Light Lean syntax highlighting. The goal is readability, not a full
% Lean-mode -- just keywords and inline UTF-8 mathematical symbols.
\lstdefinelanguage{Lean}{
  morekeywords={def,theorem,lemma,structure,where,abbrev,by,intro,exact,
                refine,rintro,obtain,rcases,subst,rw,simp,Set,Prop,Type,
                Disjoint,Walk,CDMG,Fin,True,False,not,sorry,if,then,else,
                fun,let,in},
  sensitive=true,
  morecomment=[l]{--},
  morecomment=[s]{/-}{-/},
  morestring=[b]",
}
\lstset{
  language=Lean,
  basicstyle=\ttfamily\footnotesize,
  keywordstyle=\bfseries,
  commentstyle=\itshape\color{gray!70!black},
  breaklines=true,
  showstringspaces=false,
  columns=fullflexible,
  upquote=true,
  literate=
    {α}{{$\alpha$}}1 {β}{{$\beta$}}1 {γ}{{$\gamma$}}1 {δ}{{$\delta$}}1
    {ε}{{$\varepsilon$}}1 {λ}{{$\lambda$}}1 {μ}{{$\mu$}}1 {π}{{$\pi$}}1
    {σ}{{$\sigma$}}1 {τ}{{$\tau$}}1 {Σ}{{$\Sigma$}}1 {Π}{{$\Pi$}}1
    {ℕ}{{$\mathbb{N}$}}1 {ℤ}{{$\mathbb{Z}$}}1 {ℝ}{{$\mathbb{R}$}}1
    {⊆}{{$\subseteq$}}1 {∈}{{$\in$}}1 {∉}{{$\notin$}}1 {∅}{{$\emptyset$}}1
    {∪}{{$\cup$}}1 {∩}{{$\cap$}}1 {∧}{{$\wedge$}}1 {∨}{{$\vee$}}1
    {¬}{{$\neg$}}1 {≠}{{$\neq$}}1 {≤}{{$\leq$}}1 {≥}{{$\geq$}}1
    {→}{{$\to$}}1 {←}{{$\leftarrow$}}1 {⟶}{{$\longrightarrow$}}1
    {↔}{{$\leftrightarrow$}}1 {⟨}{{$\langle$}}1 {⟩}{{$\rangle$}}1
    {⦃}{{\{\!|}}1 {⦄}{{|\!\}}}1 {×ˢ}{{$\times^{\!s}$}}1 {∃}{{$\exists$}}1
    {∀}{{$\forall$}}1 {≔}{{$:=$}}1,
}

\title{When the encoding lies:\\
how one disjointness field made a true lemma false}
\author{Documenting Progression}
\date{2026-05-29}

\begin{document}
\maketitle

\begin{abstract}
On 2026-05-28 the orchestrator marked \texttt{claim\_3\_25} (the
lecture notes' $i\sigma$-separation-under-marginalization lemma) as
solved. The artefact is in fact a \emph{disproof}: a concrete
counter-example whose Lean theorem name is
\texttt{isISigmaSeparated\_marginalize\_iff\_disproved}. Yet the
lemma is a standard, well-known result in the latent-projection
literature. This note traces, line by line, why the disproof is
nonetheless correct \emph{under our Lean encoding}, and isolates the
exact design choice -- one disjointness field on the \texttt{CDMG}
structure -- that cascaded all the way to a false lemma. The lecture
notes are not wrong; one of our encoding choices is.
\end{abstract}

\section{Where the chain starts: \texttt{def\_3\_1.CDMG}}

The lecture notes' \texttt{CDMG} definition (def 3.1 of the LN) says
\textit{``a CDMG consists of two (disjoint) sets of edges''}: a set
$E$ of directed edges (carrier $\subseteq (J \cup V) \times V$) and
a set $L$ of bidirected edges (carrier $\subseteq V \times V$,
quotiented by $(v_1, v_2) \sim (v_2, v_1)$). Because $E$ and $L$
inhabit \emph{different} ambient types in the LN's reading
($(J \cup V) \times V$ vs.\ $V \times V / \!\!\sim$), the
``disjointness'' phrase is automatically satisfied: a pair is either
a directed edge or a bidirected edge, never both -- the carriers do
not even share an object that could be in both.

Our Lean encoding (\texttt{Section3\_1/CDMG.lean}) is the verbatim
\texttt{structure} below:

\begin{lstlisting}
structure CDMG (α : Type*) where
  /-- The set of *input* nodes. -/
  J : Set α
  /-- The set of *output* nodes. -/
  V : Set α
  /-- Inputs and outputs are disjoint sets of vertices. -/
  disjoint_JV : Disjoint J V
  /-- The set of *directed* edges. -/
  E : Set (α × α)
  E_subset : E ⊆ (J ∪ V) ×ˢ V
  /-- The set of *bidirected* edges, encoded as a subset of `V × V`
  required to be symmetric and irreflexive. -/
  L : Set (α × α)
  L_subset : L ⊆ V ×ˢ V
  L_irrefl : ∀ ⦃v₁ v₂ : α⦄, (v₁, v₂) ∈ L → v₁ ≠ v₂
  L_symm   : ∀ ⦃v₁ v₂ : α⦄, (v₁, v₂) ∈ L → (v₂, v₁) ∈ L
  /-- Directed and bidirected edges are disjoint as subsets of `α × α`. -/
  disjoint_EL : Disjoint E L
\end{lstlisting}

The collapse the LN never made is right there in the carrier of $L$:
\texttt{L : Set ($\alpha$ $\times$ $\alpha$)}, with the \emph{same}
ambient type as \texttt{E}. Where the LN's $E$ and $L$ inhabit
disjoint sorts and never collide, our $E$ and $L$ are both subsets
of the same $\alpha \times \alpha$, so the same pair $(u, v)$ could
in principle live in both. To rescue the LN's ``two disjoint sets of
edges'' phrasing, the structure adds the literal-disjointness field
\texttt{disjoint\_EL : Disjoint E L} as a proof obligation.

The trade-off was deliberate: a single carrier gives downstream
proofs ergonomic uniform membership ($(u, v) \in G.E$ /
$(u, v) \in G.L$ as parallel idioms), no quotient plumbing, no
ordered-vs-unordered-pair lift gymnastics. The cost was supposed to
be one extra hypothesis. As we'll see, that cost is much bigger.

\section{The cascade reaches \texttt{def\_3\_14.marginalize}}

Marginalization (also known as latent projection) takes a CDMG $G$
and a subset $W \subseteq V$ of vertices to ``project away''. The
LN's marginalized graph $G^{\sm W}$ has vertex set
$(J, V \setminus W)$ and edge sets defined by:

\begin{itemize}
\item $(u, v) \in E^{\sm W}$ iff there is a length-$\geq 1$
      directed walk $u \tuh \cdots \tuh v$ in $G$ whose interior
      vertices (i.e.\ everything strictly between $u$ and $v$) all
      lie in $W$.
\item $(u, v) \in L^{\sm W}$ iff there is a length-$\geq 1$
      \emph{bifurcation} walk between $u$ and $v$ in $G$ with all
      interior vertices in $W$. (A bifurcation here is a walk of
      shape $u \hut \cdots \hut x \tuh \cdots \tuh v$ or
      $u \huh \cdots \huh v$ -- two arms diverging from a common
      source, possibly through a bidirected edge in the middle.)
\end{itemize}

Critically: \textbf{the LN's definition imposes no exclusion between
$E^{\sm W}$ and $L^{\sm W}$.} The same pair $(u, v)$ may inhabit
both -- e.g.\ a vertex $u$ has both a directed walk to $v$ through
$W$ \emph{and} a fork bifurcation with $v$ through $W$. Under the
LN's separate-ambient-types reading, both contributions are
recorded independently; no conflict.

Our Lean \texttt{marginalize} (\texttt{Section3\_2/Marginalization.lean})
cannot do that: it has to satisfy
\texttt{disjoint\_EL : Disjoint E L} on the output graph. So it
\emph{adds two exclusion clauses to $L^{\sm W}$}:

\begin{lstlisting}
def marginalize (G : CDMG α) (W : Set α) : CDMG α where
  J := G.J
  V := G.V \ W
  ...
  E := { p : α × α |
    p.1 ∈ G.J ∪ (G.V \ W) ∧ p.2 ∈ G.V \ W ∧
    ∃ π : Walk G p.1 p.2,
      π.IsDirected ∧ π.InteriorIn W ∧ 1 ≤ π.length }
  ...
  L := { p : α × α |
    p.1 ∈ G.V \ W ∧ p.2 ∈ G.V \ W ∧ p.1 ≠ p.2 ∧
    -- exclusion A: forbid p if a directed walk p.1 → p.2 through W exists
    (¬ ∃ π : Walk G p.1 p.2, π.IsDirected ∧ π.InteriorIn W) ∧
    -- exclusion B: forbid p if a directed walk p.2 → p.1 through W exists
    (¬ ∃ π : Walk G p.2 p.1, π.IsDirected ∧ π.InteriorIn W) ∧
    -- and require an actual bifurcation in either direction
    ((∃ π : Walk G p.1 p.2, π.IsBifurcation ∧ π.InteriorIn W) ∨
     (∃ π : Walk G p.2 p.1, π.IsBifurcation ∧ π.InteriorIn W)) }
  ...
  disjoint_EL := by
    rw [Set.disjoint_left]
    intro p hE hL
    obtain ⟨_, _, π, hπ_dir, hπ_int, _⟩ := hE
    obtain ⟨_, _, _, hnE_uv, _, _⟩ := hL
    exact hnE_uv ⟨π, hπ_dir, hπ_int⟩
\end{lstlisting}

The two exclusion clauses \texttt{(¬ $\exists$ $\pi$, ...)} are not
in the LN's definition. They are there \emph{only} because
\texttt{disjoint\_EL} on \texttt{CDMG} demands it: without them, any
pair $(u, v)$ that gets in via both a directed-walk-through-$W$ and a
bifurcation-through-$W$ would land in both $E^{\sm W}$ and $L^{\sm W}$,
violating \texttt{disjoint\_EL} on the marginalized graph.

The author of the file knew this and documented it in the design block:

\begin{quote}
``\texttt{L\^{}\{$\sm$ W\}} excludes pairs already in
\texttt{E\^{}\{$\sm$ W\}} (in either direction). This is the most
subtle deviation from the LN's literal definition, forced by the
\texttt{disjoint\_EL} field of \texttt{def\_3\_1.CDMG}. The LN treats
directed and bidirected edges as different kinds of objects \dots\
so the LN's `two (disjoint) sets of edges' is automatically satisfied
even if some pair $(u, v)$ is both a directed edge from $u$ to $v$
\emph{and} a bidirected edge between $u$ and $v$\dots\ The LN's
marginalization \emph{can} produce collisions of this kind.''
\end{quote}

The author's defense was that this deviation is benign because what
downstream proofs need is \emph{bifurcation existence}, not literal
membership in $L^{\sm W}$, and bifurcation existence survives the
exclusion (a directed walk reverses to an all-backward walk, which
counts as a bifurcation with empty right arm). Under that interpretation
\texttt{claim\_3\_16} (``marginalization preserves bifurcations'')
goes through and the encoding looks safe.

But \texttt{claim\_3\_25}'s proof needs more.

\section{The LN's lemma and the rerouting step it depends on}

The lemma in question is the LN's
\texttt{lem:stability\_separation\_marginalization}
(claim\_3\_25 in our data file):

\begin{Lem}[$i\sigma$-separation under marginalization]
Let $G = (J, V, E, L)$ be a CDMG, $A, B, C \ins J \cup V$ and
$D \ins V$ subsets such that
$D \cap (A \cup B \cup C) = \emptyset$. Then
\[ A \isPerp_G B \given C \qquad \iff \qquad
   A \isPerp_{G^{\sm D}} B \given C. \]
\end{Lem}

The LN's proof of $(\Longleftarrow)$ -- if the marginalized graph
$\sigma$-separates, so does the original -- proceeds by contrapositive.
Take a $C$-$\sigma$-open walk $\pi$ in $G$; transform it into a
$C$-$\sigma$-open walk in $G^{\sm D}$ that contradicts the
$\sigma$-separation in $G^{\sm D}$. The transformation removes
$D$-interior vertices one at a time and, at each step, may need to
\emph{reroute a fork bifurcation through a bidirected edge of
$L^{\sm D}$}. The breaking step in the LN's proof
(\texttt{graphs.tex}:1559) is exactly:

\begin{quote}
\textit{If $w \hut u \tuh b_{j+1}$ is a fork in $G$ with $u \in D$,
then in $G^{\sm \lC u \rC}$ there is a bidirected edge
$w \huh b_{j+1} \in L^{\sm \lC u \rC}$.}
\end{quote}

In the LN, this membership is immediate: the fork is a bifurcation
through $\lC u \rC$; bifurcations through $W$ deposit bidirected
edges in $L^{\sm W}$. In our encoding, the membership is conditional
on \emph{no} directed walk between $w$ and $b_{j+1}$ existing through
$\lC u \rC$ in either direction. Make one exist, and the rerouting
step the LN's proof relies on becomes unavailable.

\section{The counter-example}

The disproof of \texttt{claim\_3\_25} (in the row's tex proof file
\texttt{tex/claim\_3\_25\_proof\_ISigmaSeparation.tex}) constructs
exactly such a graph. Vertices are six labels
$v_0, b_j, u, w, b_{j+1}, v_n$ (encoded as
\texttt{Fin 6 = $\lC 0, 1, 2, 3, 4, 5\rC$} in Lean). The edges are:

\[
  E := \lC
    (v_0, b_j),\ (b_j, u),\ (u, b_{j+1}),\ (b_{j+1}, v_n),\
    (u, w),\ (w, u),\ (w, b_j)
  \rC, \qquad L := \emptyset.
\]

Read as a directed-mixed graph:

\begin{itemize}
\item A spine $v_0 \tuh b_j \tuh u \tuh b_{j+1} \tuh v_n$.
\item A 2-cycle $u \tuh w$ and $w \tuh u$ at the middle vertex.
\item An extra edge $w \tuh b_j$ closing a triangle on
      $\lC b_j, u, w \rC$.
\end{itemize}

The choices are
\[ A := \lC v_0 \rC,\quad B := \lC v_n \rC,\quad
   C := \lC b_j, w \rC,\quad D := \lC u \rC. \]
Then $D \cap (A \cup B \cup C) = \emptyset$, satisfying the lemma's
hypothesis.

\paragraph{LHS: $A \nisPerp_G B \given C$.}
The walk $\pi := \lp v_0 \tuh b_j \tuh u \tuh b_{j+1} \tuh v_n \rp$
is $C$-$\sigma$-open in $G$. Per-position:
\begin{itemize}
\item At $b_j \in C$: right-chain non-collider whose forward-out
      target $u$ lies in $\Sc^G(b_j) = \lC b_j, u, w \rC$ -- so $b_j$
      is \emph{unblockable}, and being in $C$ does not block.
\item At $u \notin C$: blockable non-collider but not in $C$, so does
      not block.
\item At $b_{j+1} \notin C$: blockable non-collider but not in $C$,
      so does not block.
\end{itemize}
No colliders (every step is forward). Hence
$A \nisPerp_G B \given C$ holds.

\paragraph{RHS: $A \isPerp_{G^{\sm \lC u \rC}} B \given C$.}
This is where the exclusion clauses bite. Marginalizing out $u$
under our \texttt{marginalize}, we compute:

\begin{itemize}
\item $E^{\sm \lC u \rC}$ contains, beside the obvious survivors,
      \emph{the pair $(w, b_{j+1})$}, via the directed walk
      $w \tuh u \tuh b_{j+1}$ (length 2, interior $\lC u \rC$).
\item $L^{\sm \lC u \rC}$ would \emph{ideally} contain
      $w \leftrightarrow b_{j+1}$, via the fork bifurcation
      $w \hut u \tuh b_{j+1}$. \textbf{Under our encoding it does
      not}: the directed walk $w \tuh u \tuh b_{j+1}$ in the
      \texttt{(u, w)} direction triggers exclusion clause A,
      forbidding $(w, b_{j+1}) \in L^{\sm \lC u \rC}$.
\end{itemize}

With $(w, b_{j+1}) \notin L^{\sm \lC u \rC}$, the LN's
rerouting move from $w$ to $b_{j+1}$ via a bidirected edge is
not available. The only path through $w$ in $G^{\sm \lC u \rC}$
uses the surviving directed edge $(w, b_{j+1}) \in E^{\sm \lC u \rC}$
as a \emph{forward step}, which makes $w$ a blockable
non-collider on any walk through it -- and $w \in C$ then blocks
the walk. Combined with $b_j \in C$ blocking (now blockable, not
unblockable, in the marginalized graph because $u$ is gone from
$\Sc$) it turns out \emph{every} walk in
$G^{\sm \lC u \rC}$ from $v_0$ to $v_n$ is $C$-$\sigma$-blocked.
So $A \isPerp_{G^{\sm \lC u \rC}} B \given C$ holds.

Combining: LHS fails, RHS holds, the equivalence is broken.

\section{The witness in Lean}

The same graph, expressed in our \texttt{CDMG (Fin 6)}:

\begin{lstlisting}
def G_witness : CDMG (Fin 6) where
  J := ∅
  V := {0, 1, 2, 3, 4, 5}
  disjoint_JV := by
    rw [Set.disjoint_left]; intro _ hJ _; exact hJ
  E := {(0, 1), (1, 2), (2, 4), (4, 5), (2, 3), (3, 2), (3, 1)}
  E_subset := by
    intro p hp
    simp only [Set.mem_insert_iff, Set.mem_singleton_iff] at hp
    rcases hp with h | h | h | h | h | h | h <;>
      (subst h; refine ⟨Or.inr ?_, ?_⟩ <;>
        simp [Set.mem_insert_iff, Set.mem_singleton_iff])
  L := ∅
  L_subset := by intro _ h; exact h.elim
  L_irrefl := by intro _ _ h; exact h.elim
  L_symm := by intro _ _ h; exact h.elim
  disjoint_EL := by
    rw [Set.disjoint_left]; intro _ _ hL; exact hL.elim
\end{lstlisting}

with the labelling $v_0 = 0,\ b_j = 1,\ u = 2,\ w = 3,\ b_{j+1} = 4,\
v_n = 5$. The conditioning set in Lean is
\texttt{C\_witness := \{1, 3\}} and the disproof theorem is:

\begin{lstlisting}
theorem isISigmaSeparated_marginalize_iff_disproved :
    ∃ (G : CDMG (Fin 6)) (A B C D : Set (Fin 6)),
      D ⊆ G.V ∧ Disjoint D (A ∪ B ∪ C) ∧
      ¬ (G.IsISigmaSeparated A B C ↔
           (G.marginalize D).IsISigmaSeparated A B C) := by
  refine ⟨G_witness, {0}, {5}, C_witness, {2}, ?_, ?_, ?_⟩
  · -- D = {2} ⊆ G_witness.V = {0, 1, 2, 3, 4, 5}.
    intro x hx
    rw [Set.mem_singleton_iff] at hx; subst hx
    exact G_witness_V_univ 2
  · -- Disjoint {2} ({0} ∪ {5} ∪ {1, 3}).
    rw [Set.disjoint_singleton_left]
    simp [C_witness, Set.mem_insert_iff, Set.mem_singleton_iff]
  · -- ¬ (LHS ↔ RHS). LHS-side: ¬ G.IsISigmaSeparated.
    -- RHS-side: marg.IsISigmaSeparated. iff forces RHS → LHS,
    -- contradicting both.
    intro hiff
    exact not_isISigmaSeparated_G_witness
      (hiff.mpr marg_isISigmaSeparated_G_witness)
\end{lstlisting}

The theorem type-checks. It is genuinely a proof of the negation of
the biconditional, instantiated on the explicit
\texttt{G\_witness}. \texttt{lake build} accepts it with no
\texttt{sorry}. It is, in every formal sense, ``proven'' -- but what
is proven is the encoding's lemma, not the LN's.

\section{Where the actual bug lives}

The chain of causes, as a directed walk through our own design space:

\begin{enumerate}
\item \texttt{def\_3\_1.CDMG} collapses directed and bidirected edges
      onto the common carrier \texttt{Set ($\alpha$ $\times$ $\alpha$)}
      and adds the field \texttt{disjoint\_EL : Disjoint E L}.
\item To keep \texttt{disjoint\_EL} valid post-marginalization,
      \texttt{def\_3\_14.marginalize} adds two ``no directed walk
      through $W$'' exclusion clauses to $L^{\sm W}$.
\item These exclusion clauses remove fork-induced bidirected edges
      whenever a directed walk also happens to exist between the
      same endpoints through $W$.
\item The LN's proof of
      \texttt{lem:stability\_separation\_marginalization} relies on
      exactly such a fork-induced bidirected edge being available
      for rerouting. With it removed, the rerouting step has no
      target, and the lemma becomes false on a structurally
      non-trivial subset of inputs.
\item \texttt{claim\_3\_25} is solved with a disproof exhibiting a
      member of that subset. Both the disproof and its Lean witness
      are \emph{correct} -- the encoding's lemma really is false.
\end{enumerate}

What is \emph{not} broken: the LN's marginalization, the LN's lemma,
the disproof itself, the formalization of the disproof in Lean, or
the workflow's machinery for recording a disprove-mode solve. What
\emph{is} broken: the encoding agreement between our
\texttt{marginalize} and the LN's. The two agree on bifurcation
existence; they disagree on literal $L^{\sm W}$ membership exactly
in the cases the lemma needs.

\section{What this means going forward}

Three actionable observations:

\paragraph{1. The fix lives upstream in \texttt{def\_3\_1}, not in
\texttt{def\_3\_14}.} As long as \texttt{CDMG} demands literal
\texttt{Disjoint E L} on a shared $\alpha \times \alpha$ carrier,
any LN-faithful marginalization will be unrepresentable. Options:
drop \texttt{disjoint\_EL} (and patch downstream call-sites that
assumed it), or move \texttt{L} to a different ambient type that
makes set-disjointness vacuous (e.g.\ \texttt{Sym2 α} via mathlib,
or a wrapped \texttt{Set ($\alpha$ $\times$ $\alpha$)} type).

\paragraph{2. The workflow's disprove pathway worked --
exactly where it shouldn't have.} The
\texttt{mistake} action was added so the manager could surrender
gracefully when the LN really is wrong. Here it surrendered
to the LN being right and the encoding being wrong, and the
orchestrator dutifully recorded a ``disproof'' of a true statement.
The verifier chain has no way to distinguish ``the LN is wrong''
from ``the encoding is wrong'' -- both look like the manager
hitting a real wall.

\paragraph{3. There is at least one more claim downstream that
inherits this issue.} The LN's $i\sigma$-separation calculus (§3.3
and beyond) uses the rerouting move repeatedly. Any LN claim whose
proof relies on a marginalization-induced bidirected edge that
also has a directed-walk version is potentially false under the
current encoding. \texttt{claim\_3\_25} is the first; it is
unlikely to be the last.

\end{document}
